\subsection{The exact row-retention budget from generalized Hamming weights}
For the same twelve sign representatives $c_1,\ldots,c_{12}$ spanning
$C$, let
\[
 E=\{(\alpha(c_1),\ldots,\alpha(c_{12})): \alpha\in C^*\}
 \subseteq\mathbb F_5^{12}.
\]
This is the ten-dimensional evaluation code used in the neighbor
census. A codimension-$h$ retained subspace $G\subset C$ has
annihilator $U\subset C^*$ of dimension $h$. Its lost old rows are
exactly the union of the supports of the evaluation words of $U$:
\[
 c_i\notin G\quad\Longleftrightarrow\quad
 \exists\alpha\in U:\alpha(c_i)\ne0.
\]
Consequently the minimum possible number of lost old sign classes
among all such $G$ is precisely the classical generalized Hamming
weight $d_h(E)$, defined as the least support size of an
$h$-dimensional subcode. This is a standard coding-theoretic resource
interpretation, not a claim to a new generalized-weight theory.
The classical reference is V. K. Wei, \emph{Generalized Hamming weights
for linear codes}, IEEE Transactions on Information Theory 37 (1991),
1412--1418, \url{https://ieeexplore.ieee.org/document/133259/}.

The two-dimensional dual has weight enumerator
$1+8z^8+8z^{10}+8z^{12}$. Its parity-check matrix $H$ has no zero
columns and projective column multiplicities $4,4,2,2,0,0$ among
the six projective points of $\mathbb F_5^2$. The final two zeros
denote absent projective directions, not zero columns. For a set
$I$ of column positions,
\[
 \dim\{e\in E:\operatorname{supp}(e)\subseteq I\}
 =|I|-\operatorname{rank}H_I,
\]
so $d_h(E)=\min\{|I|:|I|-\operatorname{rank}H_I\ge h\}$.
Any nonempty $I$ has rank at least one; rank-one sets have at most
four columns. They yield $d_1,d_2,d_3=2,3,4$. For $h\ge4$ the
rank must be two and $d_h=h+2$, attainable by choosing that many
columns. Thus
\[
 (d_1,\ldots,d_{10})=(2,3,4,6,7,8,9,10,11,12).
\]
The exact maximum retained-row counts for codimensions one through
ten are therefore $10,9,8,6,5,4,3,2,1,0$.
In particular the rank-eight, nine-sign-class benchmark used in the
completion-moment calculation has the largest possible old-row
retention for a two-dimensional annihilator. A retained space
containing five old rows can have codimension at most five.

The standard-library script \texttt{check\_retention.py} independently
checks all 4096 column subsets, writes the full witnesses to
\texttt{retention.json}, and verifies both the weight enumerator and
projective multiplicities. These are optimal retention budgets for
this specific code. They do not ensure that any corresponding
completion improves sign count or prove unknown-order existence.
